% ===== PRÉAMBULE CANONIQUE — à coller VERBATIM en tête de CHAQUE .tex =====
\documentclass[11pt,a4paper]{article}
\usepackage[utf8]{inputenc}\usepackage[T1]{fontenc}\usepackage{lmodern}
\usepackage[french]{babel}
\usepackage{amsmath,amssymb,mathtools}\usepackage{icomma}
\usepackage{siunitx}\sisetup{locale=FR}  % \qty{}{}, \si{}, \ang{}
\usepackage{xcolor}\usepackage{colortbl}
\usepackage{tikz}\usepackage{pgfplots}\pgfplotsset{compat=1.18}
\usepackage[siunitx,europeanresistors]{circuitikz} % circuits (élec.)
\usepackage[most]{tcolorbox}\usepackage{enumitem}\usepackage{array,booktabs}
\usepackage[a4paper,margin=2.2cm]{geometry}
\usetikzlibrary{arrows.meta,calc,angles,quotes,positioning,patterns,%
  backgrounds,shapes.geometric,decorations.pathmorphing,decorations.markings,intersections}
\definecolor{primary}{RGB}{222,49,99}\definecolor{primarydark}{RGB}{150,22,55}
\definecolor{defcol}{RGB}{0,114,178}\definecolor{methcol}{RGB}{52,92,148}
\definecolor{excol}{RGB}{0,135,90}\definecolor{attncol}{RGB}{200,40,55}
\tcbset{boxbase/.style={enhanced,breakable,boxrule=0.9pt,arc=2.5pt,
  left=3mm,right=3mm,top=2.3mm,bottom=2.3mm,fonttitle=\bfseries,coltitle=white}}
% --- boîtes TUTORIEL (noms EXACTS, ne pas renommer) ---
\newtcolorbox{cle}[1][]{boxbase,colback=defcol!6,colframe=defcol,
  title={\textsc{À retenir}\ifx\relax#1\relax\relax\else\textmd{ : \itshape #1}\fi}}
\newtcolorbox{methode}[1][]{boxbase,colback=methcol!7,colframe=methcol,sharp corners,
  title={\textsc{Méthode}\ifx\relax#1\relax\relax\else\textmd{ --- \itshape #1}\fi}}
\newtcolorbox{exemplebox}[1][]{boxbase,colback=excol!5,colframe=excol,
  title={\textsc{Exemple}\ifx\relax#1\relax\relax\else\textmd{ : \itshape #1}\fi}}
\newtcolorbox{piege}[1][]{boxbase,colback=attncol!6,colframe=attncol,
  title={\textsc{Piège}\ifx\relax#1\relax\relax\else\textmd{ : \itshape #1}\fi}}
\newtcolorbox{remarque}[1][]{boxbase,colback=black!4,colframe=black!45,
  title={\textsc{Remarque}\ifx\relax#1\relax\relax\else\textmd{ : \itshape #1}\fi}}
% --- boîtes EXERCICE / CORRIGÉ (noms EXACTS, auto-comptées) ---
\newtcolorbox[auto counter]{exo}[1][]{boxbase,colback=primary!4,colframe=primary,
  title={\textsc{Exercice}~\thetcbcounter\ifx\relax#1\relax\relax\else\textmd{ --- \itshape #1}\fi}}
\newtcolorbox[auto counter]{corrige}[1][]{boxbase,colback=excol!4,colframe=excol,
  title={\textsc{Corrigé}~\thetcbcounter\ifx\relax#1\relax\relax\else\textmd{ --- \itshape #1}\fi}}
\newcommand{\vv}[1]{\vec{#1}}\newcommand{\ds}{\displaystyle}
\newcommand{\R}{\mathbb{R}}\newcommand{\N}{\mathbb{N}}\newcommand{\Z}{\mathbb{Z}}
% =========================================================================
\begin{document}

\begin{corrige}[angle limite verre $\to$ air]
\[ \hat{\imath}_{\text{lim}}=\arcsin\!\left(\frac{n_2}{n_1}\right)=\arcsin\!\left(\frac{1}{1{,}5}\right)
   =\arcsin(0{,}667)\approx\ang{41.8}. \]
\end{corrige}

\begin{corrige}[angle limite eau $\to$ air]
\begin{enumerate}[label=\alph*)]
  \item $\hat{\imath}_{\text{lim}}=\arcsin\!\left(\dfrac{1}{1{,}33}\right)=\arcsin(0{,}752)\approx\ang{48.8}$.
  \item $\ang{60}>\ang{48.8}=\hat{\imath}_{\text{lim}}$ : le rayon subit une \textbf{réflexion totale},
        il ne sort pas dans l'air et reste dans l'eau.
\end{enumerate}
\end{corrige}

\begin{corrige}[réflexion totale, oui ou non ?]
\begin{enumerate}[label=\alph*)]
  \item $\hat{\imath}_{\text{lim}}=\arcsin(1/1{,}5)\approx\ang{41.8}$.
  \item On part du verre (plus réfringent) vers l'air, et $\hat{\imath}_1=\ang{50}>\ang{41.8}$ : la
        réflexion est \textbf{totale}. Le rayon reste dans le verre.
\end{enumerate}
\end{corrige}

\begin{corrige}[dans quel sens ?]
Seul le cas \textbf{(B), eau $\to$ air}, peut donner une réflexion totale : il faut partir du milieu
le \emph{plus} réfringent (l'eau) vers le moins réfringent (l'air). Dans le cas (A), air $\to$ eau
(vers un milieu plus réfringent), il y a \emph{toujours} un rayon réfracté : jamais de réflexion totale.
\end{corrige}

\begin{corrige}[angle limite verre $\to$ eau]
\begin{enumerate}[label=\alph*)]
  \item $\hat{\imath}_{\text{lim}}=\arcsin\!\left(\dfrac{n_2}{n_1}\right)=\arcsin\!\left(\dfrac{1{,}33}{1{,}5}\right)
        =\arcsin(0{,}887)\approx\ang{62.5}$.
  \item Il est \textbf{plus grand} que l'angle limite verre/air ($\ang{41.8}$). Les deux indices sont
        plus proches ($n_2/n_1$ proche de $1$), donc il faut incliner davantage le rayon pour atteindre
        la réflexion totale. Plus les milieux se ressemblent, plus l'angle limite est grand.
\end{enumerate}
\end{corrige}

\end{document}
